\documentclass[11pt,a4paper]{article}
\usepackage[utf8]{inputenc}
\usepackage[spanish]{babel}
\usepackage[T1]{fontenc}
\usepackage{geometry}
\geometry{margin=2cm}
\usepackage{xcolor}
\usepackage{listings}
\usepackage{tcolorbox}
\tcbuselibrary{most}
\usepackage{booktabs}
\usepackage{fancyhdr}
\usepackage{titlesec}
\usepackage{multicol}
\usepackage{enumitem}

\definecolor{simred}{RGB}{180,30,40}
\definecolor{codebg}{RGB}{245,245,245}
\definecolor{alertred}{RGB}{220,50,50}
\definecolor{tipgreen}{RGB}{39,174,96}
\definecolor{warnambel}{RGB}{230,126,34}
\definecolor{darkgray}{RGB}{50,50,50}
\definecolor{sqlkw}{RGB}{0,0,200}

\lstdefinestyle{python}{
  language=Python,
  backgroundcolor=\color{codebg},
  basicstyle=\ttfamily\small,
  keywordstyle=\color{simred}\bfseries,
  stringstyle=\color{teal},
  commentstyle=\color{gray}\itshape,
  numbers=left, numberstyle=\tiny\color{gray},
  frame=single, framerule=0.4pt, rulecolor=\color{simred!40},
  breaklines=true, tabsize=2, showstringspaces=false
}

\lstdefinelanguage{SQL2}{
  keywords={SELECT,FROM,WHERE,INSERT,INTO,VALUES,UPDATE,SET,DELETE,CREATE,TABLE,
            JOIN,LEFT,GROUP,BY,ORDER,HAVING,LIMIT,WITH,AS,ON,DISTINCT,RETURNING,
            COUNT,SUM,AVG,NOT,NULL,AND,OR,INNER,OUTER},
  sensitive=false,
  morestring=[b]', morecomment=[l]{--}
}

\lstdefinestyle{sql}{
  language=SQL2,
  backgroundcolor=\color{codebg},
  basicstyle=\ttfamily\small,
  keywordstyle=\color{sqlkw}\bfseries,
  stringstyle=\color{teal},
  commentstyle=\color{gray}\itshape,
  numbers=left, numberstyle=\tiny\color{gray},
  frame=single, framerule=0.4pt, rulecolor=\color{sqlkw!20},
  breaklines=true, tabsize=2, showstringspaces=false
}

\newtcolorbox{preguntabox}[1]{
  colback=simred!5, colframe=simred,
  title={\textbf{#1}}, fonttitle=\small\bfseries,
  before upper={\setlength{\parskip}{4pt}}
}
\newtcolorbox{respuestabox}{
  colback=tipgreen!8, colframe=tipgreen!60,
  title={\textbf{Respuesta ideal}}, fonttitle=\small\bfseries
}
\newtcolorbox{tipbox}{colback=tipgreen!10,colframe=tipgreen,title=\textbf{Tip estratégico},fonttitle=\small\bfseries}
\newtcolorbox{alertbox}{colback=alertred!10,colframe=alertred,title=\textbf{Trampa común},fonttitle=\small\bfseries}

\pagestyle{fancy}
\fancyhf{}
\lhead{\textcolor{simred}{\textbf{Simulacros de Entrevista}}}
\rhead{\textcolor{gray}{Preparación Técnica — Toku}}
\cfoot{\thepage}

\titleformat{\section}{\large\bfseries\color{simred}}{}{0em}{}[\titlerule]
\titleformat{\subsection}{\normalsize\bfseries\color{darkgray}}{}{0em}{}

\begin{document}

\begin{titlepage}
  \centering
  \vspace*{2cm}
  {\Huge\bfseries\color{simred} Simulacros de Entrevista\par}
  \vspace{0.5cm}
  {\large\color{gray} Toku — Software Engineer 2026\par}
  \vspace{0.3cm}
  \textcolor{simred!40}{\rule{8cm}{2pt}}
  \vspace{1.5cm}

  \begin{tcolorbox}[colback=simred!5,colframe=simred,width=11cm]
    \centering
    \textbf{Formato real de entrevistas en startups chilenas tech:}\\[6pt]
    30 min cultura + motivación\\
    45-60 min técnico (código en vivo o take-home)\\
    15 min preguntas del candidato
  \end{tcolorbox}

  \vspace{1.5cm}
  \begin{tabular}{ll}
    \textcolor{simred}{\textbullet} & Preguntas de diseño de sistemas\\
    \textcolor{simred}{\textbullet} & Coding challenges de pagos\\
    \textcolor{simred}{\textbullet} & Preguntas de cultura y trade-offs\\
    \textcolor{simred}{\textbullet} & SQL en vivo\\
    \textcolor{simred}{\textbullet} & Tus preguntas para ellos (estratégicas)\\
  \end{tabular}
  \vfill
  {\small\color{gray} Marzo 2026 — Estudio Personal\par}
\end{titlepage}

\tableofcontents
\newpage

% ─── BLOQUE 1: DISEÑO DE SISTEMAS ────────────────────────────────────────────
\section{Diseño de Sistemas}

\begin{preguntabox}{Pregunta 1 — Sistema de reintentos de cobro}
\textit{"Tienes que diseñar un sistema que reintente automáticamente los cobros fallidos.
Un cobro puede fallar por fondos insuficientes (reintentable) o por cuenta cerrada (no reintentable).
Describe cómo lo diseñarías."}
\end{preguntabox}

\begin{respuestabox}
\textbf{Estructura de datos:}
\begin{lstlisting}[style=sql,numbers=none]
CREATE TABLE retry_queue (
  id UUID PRIMARY KEY DEFAULT gen_random_uuid(),
  invoice_id UUID NOT NULL REFERENCES invoices(id),
  intentos INTEGER DEFAULT 0,
  max_intentos INTEGER DEFAULT 5,
  proximo_intento TIMESTAMPTZ NOT NULL,
  error_codigo TEXT,
  estado TEXT DEFAULT 'pendiente'
    CHECK (estado IN ('pendiente','procesando','completado','abandonado'))
);
\end{lstlisting}

\textbf{Lógica de reintentos:}
\begin{lstlisting}[style=python,numbers=none]
ERRORES_NO_RETRIABLES = {"CUENTA_CERRADA", "MANDATO_REVOCADO", "FRAUDE"}

async def procesar_fallo(invoice_id, error_codigo):
    if error_codigo in ERRORES_NO_RETRIABLES:
        await marcar_como_fallido_permanente(invoice_id)
        await notificar_organizacion(invoice_id, "no_retriable")
        return

    entry = await obtener_retry_entry(invoice_id)
    if entry["intentos"] >= entry["max_intentos"]:
        await abandonar(invoice_id)
        return

    # Backoff exponencial: 1h, 4h, 24h, 48h, 72h
    horas = [1, 4, 24, 48, 72][entry["intentos"]]
    proximo = NOW() + INTERVAL f'{horas} hours'

    await actualizar_retry_entry(invoice_id, entry["intentos"]+1, proximo)
\end{lstlisting}

\textbf{Worker (Celery/cron job):}
\begin{lstlisting}[style=python,numbers=none]
# Cada 15 minutos: tomar tareas listas para procesar
SELECT * FROM retry_queue
WHERE estado = 'pendiente'
  AND proximo_intento <= NOW()
  AND intentos < max_intentos
FOR UPDATE SKIP LOCKED
LIMIT 100;
\end{lstlisting}
\end{respuestabox}

\begin{tipbox}
Los entrevistadores buscan: (1) separación de errores retriables vs no retriables,
(2) backoff exponencial, (3) concurrencia sin doble procesamiento (SKIP LOCKED),
(4) notificación al fallar definitivamente.
\end{tipbox}

\vspace{8pt}

\begin{preguntabox}{Pregunta 2 — Diseña el endpoint POST /transactions}
\textit{"Diseña el endpoint que inicia un cobro. Considera concurrencia, idempotencia y errores."}
\end{preguntabox}

\begin{respuestabox}
\begin{lstlisting}[style=python,numbers=none]
@router.post("/transactions", status_code=201)
async def crear_transaccion(
    payload: TransactionCreate,
    idempotency_key: str = Header(...),  # obligatorio
    org: dict = Depends(get_current_org),
    db = Depends(get_db),
):
    # 1. Idempotencia: verificar si ya existe con esta key
    existente = await db.fetchrow(
        "SELECT * FROM transactions WHERE idempotency_key = $1", idempotency_key
    )
    if existente:
        return dict(existente)  # mismo resultado

    # 2. Verificar que la invoice existe y es 'pending', LOCK para evitar doble cobro
    async with db.transaction():
        invoice = await db.fetchrow(
            "SELECT * FROM invoices WHERE id=$1 AND estado='pending' FOR UPDATE",
            payload.invoice_id
        )
        if not invoice:
            raise HTTPException(404, "Invoice no encontrada o ya procesada")

        # 3. Insertar la transaccion
        txn = await db.fetchrow(
            """INSERT INTO transactions (invoice_id, monto, idempotency_key, estado)
               VALUES ($1,$2,$3,'pending') RETURNING *""",
            payload.invoice_id, invoice["monto"], idempotency_key
        )

    # 4. Enviar al sistema bancario de forma asincrona
    background_tasks.add_task(procesar_cobro_bancario, str(txn["id"]))

    return dict(txn)
\end{lstlisting}
\end{respuestabox}

% ─── BLOQUE 2: CODING CHALLENGES ─────────────────────────────────────────────
\section{Coding Challenges}

\begin{preguntabox}{Challenge 1 — Procesar CSV de clientes (estilo SFTP de Toku)}
\textit{Escribe una función que lea un CSV de clientes y los cree en Toku vía API.
Maneja errores por fila sin abortar todo el proceso. Devuelve un reporte.}
\end{preguntabox}

\begin{respuestabox}
\begin{lstlisting}[style=python]
import csv
import io
from dataclasses import dataclass, field
from typing import List

@dataclass
class ResultadoImport:
    exitosos: int = 0
    fallidos: int = 0
    errores: List[dict] = field(default_factory=list)

async def importar_clientes_csv(
    csv_content: str,
    toku_client,
) -> ResultadoImport:
    resultado = ResultadoImport()
    reader = csv.DictReader(io.StringIO(csv_content))

    for num_fila, fila in enumerate(reader, start=2):  # start=2 por header
        try:
            # Validacion basica
            if not fila.get("email") or not fila.get("nombre"):
                raise ValueError("email y nombre son obligatorios")

            await toku_client.create_customer({
                "external_id": fila.get("id"),
                "nombre": fila["nombre"].strip(),
                "email": fila["email"].strip().lower(),
                "telefono": fila.get("telefono", "").strip() or None,
            })
            resultado.exitosos += 1

        except Exception as e:
            resultado.fallidos += 1
            resultado.errores.append({
                "fila": num_fila,
                "datos": fila,
                "error": str(e)
            })

    return resultado
\end{lstlisting}
\end{respuestabox}

\vspace{8pt}

\begin{preguntabox}{Challenge 2 — Calcular conciliación}
\textit{Dada una lista de transacciones de tu sistema y una de Toku,
encuentra las discrepancias (en tu sistema pero no en Toku, y viceversa).}
\end{preguntabox}

\begin{respuestabox}
\begin{lstlisting}[style=python]
from dataclasses import dataclass
from typing import List, Dict

@dataclass
class Discrepancia:
    tipo: str  # 'solo_en_tu_sistema' | 'solo_en_toku' | 'monto_diferente'
    transaction_id: str
    detalle: str

def conciliar_transacciones(
    transacciones_tuyas: List[Dict],
    transacciones_toku: List[Dict],
) -> List[Discrepancia]:
    # Indexar por ID para O(n) en lugar de O(n^2)
    idx_tuyas = {t["id"]: t for t in transacciones_tuyas}
    idx_toku = {t["id"]: t for t in transacciones_toku}

    discrepancias = []

    # Transacciones en tu sistema pero no en Toku
    for tid, txn in idx_tuyas.items():
        if tid not in idx_toku:
            discrepancias.append(Discrepancia(
                tipo="solo_en_tu_sistema",
                transaction_id=tid,
                detalle=f"Monto: {txn['monto']}"
            ))
        elif abs(txn["monto"] - idx_toku[tid]["monto"]) > 0:
            discrepancias.append(Discrepancia(
                tipo="monto_diferente",
                transaction_id=tid,
                detalle=f"Tu sistema: {txn['monto']}, Toku: {idx_toku[tid]['monto']}"
            ))

    # Transacciones en Toku pero no en tu sistema
    for tid in idx_toku:
        if tid not in idx_tuyas:
            discrepancias.append(Discrepancia(
                tipo="solo_en_toku",
                transaction_id=tid,
                detalle=f"Monto: {idx_toku[tid]['monto']}"
            ))

    return discrepancias
\end{lstlisting}
\end{respuestabox}

% ─── BLOQUE 3: SQL EN VIVO ───────────────────────────────────────────────────
\section{SQL en Vivo}

\begin{preguntabox}{SQL 1 — Tasa de éxito por método de pago}
\textit{"¿Qué método de pago tiene mejor tasa de éxito en el último trimestre?"}
\end{preguntabox}

\begin{respuestabox}
\begin{lstlisting}[style=sql]
SELECT
  pm.tipo AS metodo_pago,
  COUNT(t.id) AS total_intentos,
  COUNT(t.id) FILTER (WHERE t.estado = 'completed') AS exitosos,
  ROUND(
    100.0 * COUNT(t.id) FILTER (WHERE t.estado = 'completed')
    / NULLIF(COUNT(t.id), 0), 2
  ) AS tasa_exito_pct
FROM transactions t
JOIN payment_methods pm ON pm.id = t.payment_method_id
WHERE t.created_at >= NOW() - INTERVAL '3 months'
GROUP BY pm.tipo
ORDER BY tasa_exito_pct DESC;
\end{lstlisting}
\end{respuestabox}

\vspace{6pt}

\begin{preguntabox}{SQL 2 — Clientes que nunca han pagado}
\textit{"Lista los clientes que tienen facturas pero nunca han completado un pago."}
\end{preguntabox}

\begin{respuestabox}
\begin{lstlisting}[style=sql]
SELECT c.id, c.nombre, c.email, COUNT(i.id) AS facturas_totales
FROM customers c
JOIN invoices i ON i.customer_id = c.id
WHERE NOT EXISTS (
  SELECT 1 FROM invoices i2
  JOIN transactions t ON t.invoice_id = i2.id
  WHERE i2.customer_id = c.id
    AND t.estado = 'completed'
)
GROUP BY c.id, c.nombre, c.email
ORDER BY facturas_totales DESC;
\end{lstlisting}
\end{respuestabox}

% ─── BLOQUE 4: TRADE-OFFS ────────────────────────────────────────────────────
\section{Preguntas de Trade-offs y Cultura}

\begin{preguntabox}{Trade-off 1 — REST vs Webhooks para notificaciones}
\textit{"Un cliente te pide que en lugar de webhooks, pueda hacer polling para saber si
su factura fue pagada. ¿Qué le dices?"}
\end{preguntabox}

\begin{respuestabox}
El polling es técnicamente válido pero tiene desventajas claras en este contexto:

\begin{itemize}[leftmargin=*]
  \item \textbf{Latencia}: el cliente pregunta cada X segundos; el webhook notifica en segundos
  \item \textbf{Carga}: 1000 clientes haciendo polling = 1000 requests/intervalo aunque no haya cambios
  \item \textbf{Complejidad}: el cliente necesita implementar el loop, manejar backoff, etc.
  \item \textbf{Costo}: más requests = más cómputo = más costos para ambos lados
\end{itemize}

\textbf{Cuándo el polling SÍ tiene sentido:}
\begin{itemize}[leftmargin=*]
  \item El cliente no puede exponer un servidor público (detrás de un firewall)
  \item Solo necesita verificar el estado una vez (no eventos en tiempo real)
  \item El intervalo de verificación es largo (ej: revisión diaria)
\end{itemize}

\textbf{Respuesta ideal}: implementar ambos. El endpoint \texttt{GET /invoices/\{id\}} siempre
disponible para polling ocasional; webhooks para notificaciones en tiempo real.
\end{respuestabox}

\vspace{8pt}

\begin{preguntabox}{Trade-off 2 — Microservicio vs monolito}
\textit{"¿Cuándo dividirías un microservicio en dos? ¿Cuál es el costo?"}
\end{preguntabox}

\begin{respuestabox}
\textbf{Dividir cuando:}
\begin{itemize}[leftmargin=*]
  \item Escalado diferente: el procesador de cobros necesita más CPU que el CRUD de clientes
  \item Equipos distintos con deploys independientes
  \item Dominios claramente separados con contratos estables
  \item Un componente falla frecuentemente y afecta al resto
\end{itemize}

\textbf{El costo de dividir:}
\begin{itemize}[leftmargin=*]
  \item Latencia de red entre servicios (antes era una llamada de función)
  \item Transacciones distribuidas (ya no puedes hacer BEGIN/COMMIT en una BD)
  \item Más infraestructura: CI/CD, monitoreo, service discovery
  \item Complejidad operacional: debugging requiere correlación de logs entre servicios
\end{itemize}

\textbf{Para Toku en etapa actual}: probablemente empezar con un monolito bien modularizado
y extraer servicios solo cuando hay un problema concreto de escala o team ownership.
\end{respuestabox}

% ─── BLOQUE 5: TUS PREGUNTAS ─────────────────────────────────────────────────
\section{Tus Preguntas para Ellos (Estratégicas)}

\begin{tcolorbox}[colback=simred!5,colframe=simred,title=\textbf{Por qué importa hacer buenas preguntas}]
Las preguntas que haces revelan cómo piensas. En Toku especialmente valoran gente
que cuestiona y tiene criterio. No preguntes cosas que puedes googlear.
\end{tcolorbox}

\subsection{Preguntas sobre el stack y producto}

\begin{enumerate}[leftmargin=*]
  \item \textit{"La oferta menciona Python/FastAPI y TypeScript/Vue.js — ¿hay servicios en producción
  donde mezclan ambos en el mismo equipo o están divididos por squad?"}

  \item \textit{"Vi que en 2022 tenían Firestore en el stack y ya no aparece. ¿Qué pasó con eso?
  ¿Fue una exploración que descartaron o migraron?"}

  \item \textit{"¿Cómo manejan las diferencias regulatorias entre Chile, México y Brasil?
  ¿Un servicio por país o un sistema genérico con configuración?"}

  \item \textit{"¿Qué pasa en producción cuando el sistema bancario está caído o lento?
  ¿Tienen circuit breakers implementados?"}
\end{enumerate}

\subsection{Preguntas sobre el equipo y trabajo}

\begin{enumerate}[leftmargin=*]
  \item \textit{"¿Cómo está estructurado el equipo de engineering? ¿Hay squads por producto
  o por capa técnica (frontend, backend, infra)?"}

  \item \textit{"¿Cuál fue el incidente de producción más difícil que tuvieron el último año
  y qué aprendieron de él?"}

  \item \textit{"¿Cómo es el proceso de code review? ¿Hay un estándar de cuánto tiempo
  puede estar un PR esperando revisión?"}

  \item \textit{"La oferta dice que buscan para el rol de Integrations Engineer alguien
  que quiera crecer hacia liderazgo técnico. ¿Cómo se ve ese camino en Toku concretamente?"}
\end{enumerate}

\subsection{Pregunta de cierre (poderosa)}

\begin{tcolorbox}[colback=tipgreen!8,colframe=tipgreen]
\textit{"¿Cuál es el mayor desafío técnico no resuelto que tiene el equipo hoy?
¿En qué área un nuevo SE podría tener mayor impacto en los primeros 6 meses?"}

Esta pregunta demuestra: orientación a impacto, que piensas en la empresa (no solo en ti),
y te da información real sobre qué van a pedirte hacer.
\end{tcolorbox}

% ─── RESUMEN FINAL ───────────────────────────────────────────────────────────
\newpage
\section{Guía Rápida: Día de la Entrevista}

\begin{multicols}{2}

\textbf{\textcolor{simred}{Los 5 conceptos más preguntados}}
\begin{enumerate}[leftmargin=*,noitemsep]
  \item Idempotencia (siempre)
  \item SELECT FOR UPDATE / bloqueo de filas
  \item Backoff exponencial en reintentos
  \item Verificación HMAC en webhooks
  \item Trade-off: cuándo usar cada tecnología
\end{enumerate}

\textbf{\textcolor{simred}{Frases que resuenan en Toku}}
\begin{itemize}[leftmargin=*,noitemsep]
  \item "El trade-off que veo es..."
  \item "Para garantizar idempotencia haría..."
  \item "En producción esto podría fallar si..."
  \item "Dependería del volumen / la frecuencia / el SLA"
  \item "Empezaría simple y lo extraería cuando..."
\end{itemize}

\columnbreak

\textbf{\textcolor{simred}{Lo que NO hacer}}
\begin{itemize}[leftmargin=*,noitemsep]
  \item Dar la respuesta sin razonar el camino
  \item Decir "siempre uso microservicios" (oversimplificado)
  \item No preguntar nada al final
  \item Ignorar los edge cases (doble cobro, timeout, red caída)
  \item Decir "no sé" sin intentar razonar
\end{itemize}

\textbf{\textcolor{simred}{Si no sabes algo}}
\begin{tcolorbox}[colback=tipgreen!8,colframe=tipgreen]
"No lo tengo claro, pero lo abordaría así:
[razonar en voz alta]. ¿Eso es lo que buscan
o hay algo que me estoy perdiendo?"
\end{tcolorbox}

\end{multicols}

\end{document}
